
\section{Sentence Objects}\fbeg{}\begin{frame}{\ft{Sentence Objects}}

\vspace{-.5em}	

%\MyDiamond{}
{\Large\fontfamily{uhv}\selectfont
%\begin{center}
\hspace*{2pt}\begin{minipage}{.96\textwidth}
\vspace{4pt}

%\definecolor{blback}{RGB}{0,100,100}	
%\definecolor{blfront}{RGB}{0,100,50}

%\fcolorbox{lqboutercolor}{lqbinnercolor}{\begin{minipage}{\textwidth}%
		
\begin{lightquadblockc}{0,0.1,0.1}{\parbox{21cm}{\centering \vspace{10pt}SENTENCES ARE THE BASIC UNITS OF DISCOURSE\vspace{3pt}}}
\begin{center}\begin{minipage}{\textwidth}
{\LARGE \setlength{\leftmargini}{30pt}\begin{enumerate}
\dmitem \textbf{Written Natural Language} \hspace{.5em} XML, JSON, or other markup is merely 
an approximate serialization of textual and discursive structures

{\vsp}

\dmitem \textbf{Sentence Nesting/Divergence Points}  \hspace{.5em} Sentences usually have a simple sequential 
organization, but not always (cf. footnotes, block quotes, enumerations)

{\vsp}

\dmitem \textbf{Object-Oriented Implementations}  \hspace{.5em} Textual objects should be 
conceptualized and, as much as possible, exposed to computer code as class-instances 
for Object-Oriented languages such as Java or C++.  Markup should be intended 
for serialization, not to provide the definitive structure of documents for 
querying and other computational requirements.

{\vsp}

\dmitem \textbf{Multiple Scales}  \hspace{.5em} Keyphrases and annotations provide a 
sub-sentence level of organization.  Paragraphs and (sub)sections provide a 
super-sentence level.
\end{enumerate}
}\end{minipage}
\end{center}
\end{lightquadblockc}
\end{minipage}

%}
%\end{minipage}
%\end{center}
}

\end{frame}

